\section{Invertible sheaves on Proj}
\label{section-invertible-on-proj}

\noindent
Recall from Algebra, Section \ref{algebra-section-graded}
the construction of the twisted module $M(n)$ associated
to a graded module over a graded ring.

\begin{definition}
\label{definition-twist}
Let $S$ be a graded ring. Let $X = \text{Proj}(S)$.
\begin{enumerate}
\item We define $\mathcal{O}_X(n) = \widetilde{S(n)}$.
This is called the $n$th
{\it twist of the structure sheaf of $\text{Proj}(S)$}.
\item For any sheaf of $\mathcal{O}_X$-modules $\mathcal{F}$ we set
$\mathcal{F}(n) = \mathcal{F} \otimes_{\mathcal{O}_X} \mathcal{O}_X(n)$.
\end{enumerate}
\end{definition}

\noindent
We are going to use Lemma \ref{lemma-widetilde-tensor}
to construct some canonical maps.
Since $S(n) \otimes_S S(m) = S(n + m)$ we see that there
are canonical maps
\begin{equation}
\label{equation-multiply}
\mathcal{O}_X(n) \otimes_{\mathcal{O}_X} \mathcal{O}_X(m)
\longrightarrow
\mathcal{O}_X(n + m).
\end{equation}
These maps are not isomorphisms in general, see the example in
Remark \ref{remark-not-isomorphism}. The same example shows
that $\mathcal{O}_X(n)$ is {\it not} an invertible sheaf on $X$ in
general. Tensoring with an arbitrary $\mathcal{O}_X$-module $\mathcal{F}$
we get maps
\begin{equation}
\label{equation-multiply-on-sheaf}
\mathcal{O}_X(n) \otimes_{\mathcal{O}_X} \mathcal{F}(m)
\longrightarrow
\mathcal{F}(n + m).
\end{equation}
The maps (\ref{equation-map-global-sections-degree-n})
on global sections give a map of graded rings
\begin{equation}
\label{equation-global-sections}
S \longrightarrow \bigoplus\nolimits_{n \geq 0} \Gamma(X, \mathcal{O}_X(n)).
\end{equation}
And for an arbitrary $\mathcal{O}_X$-module $\mathcal{F}$ the maps
(\ref{equation-multiply-on-sheaf}) give a graded module structure
\begin{equation}
\label{equation-global-sections-module}
\bigoplus\nolimits_{n \geq 0} \Gamma(X, \mathcal{O}_X(n))
\times
\bigoplus\nolimits_{m \in \mathbf{Z}} \Gamma(X, \mathcal{F}(m))
\longrightarrow
\bigoplus\nolimits_{m \in \mathbf{Z}} \Gamma(X, \mathcal{F}(m))
\end{equation}
and via (\ref{equation-global-sections}) also a $S$-module structure.
More generally, given any graded $S$-module
$M$ we have $M(n) = M \otimes_S S(n)$. Hence we get maps
\begin{equation}
\label{equation-multiply-more-generally}
\widetilde M(n)
=
\widetilde M
\otimes_{\mathcal{O}_X}
\mathcal{O}_X(n)
\longrightarrow
\widetilde{M(n)}.
\end{equation}
On global sections (\ref{equation-map-global-sections-degree-n})
defines a map of graded $S$-modules
\begin{equation}
\label{equation-global-sections-more-generally}
M \longrightarrow
\bigoplus\nolimits_{n \in \mathbf{Z}} \Gamma(X, \widetilde{M(n)}).
\end{equation}
Here is an important fact which follows basically immediately from the
definitions.

\begin{lemma}
\label{lemma-when-invertible}
Let $S$ be a graded ring. Set $X = \text{Proj}(S)$.
Let $f \in S$ be homogeneous of degree $d > 0$.
The sheaves $\mathcal{O}_X(nd)|_{D_{+}(f)}$ are invertible,
and in fact trivial for all $n \in \mathbf{Z}$
(see Modules, Definition \ref{modules-definition-invertible}).
The maps (\ref{equation-multiply}) restricted to $D_{+}(f)$
$$
\mathcal{O}_X(nd)|_{D_{+}(f)} \otimes_{\mathcal{O}_{D_{+}(f)}}
\mathcal{O}_X(m)|_{D_{+}(f)}
\longrightarrow
\mathcal{O}_X(nd + m)|_{D_{+}(f)},
$$
the maps (\ref{equation-multiply-on-sheaf}) restricted to $D_+(f)$
$$
\mathcal{O}_X(nd)|_{D_{+}(f)} \otimes_{\mathcal{O}_{D_{+}(f)}}
\mathcal{F}(m)|_{D_{+}(f)}
\longrightarrow
\mathcal{F}(nd + m)|_{D_{+}(f)},
$$
and the maps (\ref{equation-multiply-more-generally})
restricted to $D_{+}(f)$
$$
\widetilde M(nd)|_{D_{+}(f)}
=
\widetilde M|_{D_{+}(f)}
\otimes_{\mathcal{O}_{D_{+}(f)}}
\mathcal{O}_X(nd)|_{D_{+}(f)}
\longrightarrow
\widetilde{M(nd)}|_{D_{+}(f)}
$$
are isomorphisms for all $n, m \in \mathbf{Z}$.
\end{lemma}

\begin{proof}
The graded $S$-module maps $S \to S(nd)$, and $M \to M(nd)$, given by
$x \mapsto f^n x$ become isomorphisms after inverting $f$. The first shows that
$S_{(f)} \cong S(nd)_{(f)}$ which gives an isomorphism
$\mathcal{O}_{D_{+}(f)} \cong \mathcal{O}_X(nd)|_{D_{+}(f)}$.
The second shows that the map
$S(nd)_{(f)} \otimes_{S_{(f)}} M_{(f)} \to M(nd)_{(f)}$
is an isomorphism. The case of the map (\ref{equation-multiply-on-sheaf})
is a consequence of the case of the map (\ref{equation-multiply}).
\end{proof}

\begin{lemma}
\label{lemma-apply-modules}
Let $S$ be a graded ring. Let $M$ be a graded $S$-module.
Set $X = \text{Proj}(S)$. Assume $X$ is covered by the standard
opens $D_+(f)$ with $f \in S_1$, e.g., if $S$ is generated by $S_1$
over $S_0$. Then the sheaves $\mathcal{O}_X(n)$
are invertible and the maps
(\ref{equation-multiply}), (\ref{equation-multiply-on-sheaf}), and
(\ref{equation-multiply-more-generally}) are isomorphisms.
In particular, these maps induce isomorphisms
$$
\mathcal{O}_X(1)^{\otimes n} \cong
\mathcal{O}_X(n)
\quad
\text{and}
\quad
\widetilde{M} \otimes_{\mathcal{O}_X} \mathcal{O}_X(n) =
\widetilde{M}(n) \cong \widetilde{M(n)}
$$
Thus (\ref{equation-map-global-sections-degree-n}) becomes a map
\begin{equation}
\label{equation-map-global-sections-degree-n-simplified}
M_n \longrightarrow \Gamma(X, \widetilde{M}(n))
\end{equation}
and (\ref{equation-global-sections-more-generally}) becomes a map
\begin{equation}
\label{equation-global-sections-more-generally-simplified}
M \longrightarrow
\bigoplus\nolimits_{n \in \mathbf{Z}} \Gamma(X, \widetilde{M}(n)).
\end{equation}
\end{lemma}

\begin{proof}
Under the assumptions of the lemma $X$ is covered by the
open subsets $D_{+}(f)$ with $f \in S_1$ and the
lemma is a consequence of Lemma \ref{lemma-when-invertible} above.
\end{proof}

\begin{lemma}
\label{lemma-where-invertible}
Let $S$ be a graded ring. Set $X = \text{Proj}(S)$. Fix $d \geq 1$ an
integer. The following open subsets of $X$ are equal:
\begin{enumerate}
\item The largest open subset $W = W_d \subset X$ such that
each $\mathcal{O}_X(dn)|_W$ is invertible and all the
multiplication maps
$\mathcal{O}_X(nd)|_W \otimes_{\mathcal{O}_W} \mathcal{O}_X(md)|_W
\to \mathcal{O}_X(nd + md)|_W$
(see \ref{equation-multiply}) are isomorphisms.
\item The union of the open subsets $D_{+}(fg)$ with
$f, g \in S$ homogeneous and $\deg(f) = \deg(g) + d$.
\end{enumerate}
Moreover, all the maps
$\widetilde M(nd)|_W = \widetilde M|_W \otimes_{\mathcal{O}_W}
\mathcal{O}_X(nd)|_W \to \widetilde{M(nd)}|_W$
(see \ref{equation-multiply-more-generally}) are isomorphisms.
\end{lemma}

\begin{proof}
If $x \in D_{+}(fg)$ with $\deg(f) = \deg(g) + d$ then
on $D_{+}(fg)$ the sheaves $\mathcal{O}_X(dn)$
are generated by the element $(f/g)^n = f^{2n}/(fg)^n$. This implies $x$
is in the open subset $W$ defined in (1) by arguing as in the
proof of Lemma \ref{lemma-when-invertible}.

\medskip\noindent
Conversely, suppose that $\mathcal{O}_X(dn)$ is free of rank 1
for all $n$ in an open neighbourhood $V$ of $x \in X$ and all the
multiplication maps
$\mathcal{O}_X(nd)|_V \otimes_{\mathcal{O}_V} \mathcal{O}_X(md)|_V
\to \mathcal{O}_X(nd + md)|_V$ are isomorphisms.
We may choose $h \in S_{+}$ homogeneous such that $x \in D_{+}(h) \subset V$.
By the definition of the twists of the structure sheaf we conclude there
exists an element $s$ of $(S_h)_d$ such that $s^n$ is a basis of $(S_h)_{nd}$
as a module over $S_{(h)}$ for all $n \in \mathbf{Z}$.
We may write
$s = f/h^m$ for some $m \geq 1$ and $f \in S_{d + m \deg(h)}$.
Set $g = h^m$ so $s = f/g$. Note that $x \in D_{+}(g)$ by construction.
Note that $g^d \in (S_h)_{d\deg(g)}$.
By assumption we can write this as a multiple of
$s^{\deg(g)} = f^{\deg(g)}/g^{\deg(g)}$, say
$g^d = a/g^e \cdot f^{\deg(g)}/g^{\deg(g)}$.
Then we conclude that $g^{d + e + \deg(g)} = a f^{\deg(g)}$
and hence also $x \in D_{+}(f)$. So $x$ is an element of the set defined
in (2).

\medskip\noindent
The existence of the generating section $s = f/g$ over
the affine open $D_{+}(fg)$ whose
powers freely generate the sheaves of modules
$\mathcal{O}_X(nd)$ easily implies that the multiplication maps
$\widetilde M(nd)|_W = \widetilde M|_W \otimes_{\mathcal{O}_W}
\mathcal{O}_X(nd)|_W \to \widetilde{M(nd)}|_W$
(see \ref{equation-multiply-more-generally})
are isomorphisms. Compare with the proof of Lemma \ref{lemma-when-invertible}.
\end{proof}

\noindent
Recall from Modules, Lemma \ref{modules-lemma-s-open}
that given an invertible sheaf $\mathcal{L}$ on a locally ringed
space $X$, and given a global section $s$ of $\mathcal{L}$
the set $X_s = \{x \in X \mid s \not \in \mathfrak m_x\mathcal{L}_x\}$
is open.

\begin{lemma}
\label{lemma-principal-open}
Let $S$ be a graded ring. Set $X = \text{Proj}(S)$. Fix $d \geq 1$ an
integer. Let $W = W_d \subset X$ be the open subscheme defined in
Lemma \ref{lemma-where-invertible}. Let $n \geq 1$ and $f \in S_{nd}$.
Denote $s \in \Gamma(W, \mathcal{O}_W(nd))$ the section which is
the image of $f$ via (\ref{equation-global-sections}) restricted to $W$. Then
$$
W_s = D_{+}(f) \cap W.
$$
\end{lemma}

\begin{proof}
Let $D_{+}(ab) \subset W$ be a standard affine open with
$a, b \in S$ homogeneous and $\deg(a) = \deg(b) + d$.
Note that $D_{+}(ab) \cap D_{+}(f) = D_{+}(abf)$.
On the other hand the restriction of $s$ to $D_{+}(ab)$
corresponds to the element $f/1 = b^nf/a^n (a/b)^n \in (S_{ab})_{nd}$.
We have seen in the proof of Lemma \ref{lemma-where-invertible} that
$(a/b)^n$ is a generator for $\mathcal{O}_W(nd)$ over $D_{+}(ab)$.
We conclude that $W_s \cap D_{+}(ab)$ is the principal open
associated to $b^nf/a^n \in \mathcal{O}_X(D_{+}(ab))$.
Thus the result of the lemma is clear.
\end{proof}

\noindent
The following lemma states the properties that we will later use to
characterize schemes with an ample invertible sheaf.

\begin{lemma}
\label{lemma-ample-on-proj}
Let $S$ be a graded ring.
Let $X = \text{Proj}(S)$.
Let $Y \subset X$ be a quasi-compact open subscheme.
Denote $\mathcal{O}_Y(n)$ the restriction of
$\mathcal{O}_X(n)$ to $Y$.
There exists an integer $d \geq 1$ such that
\begin{enumerate}
\item the subscheme $Y$ is contained in the open $W_d$ defined
in Lemma \ref{lemma-where-invertible},
\item the sheaf $\mathcal{O}_Y(dn)$ is invertible for all $n \in \mathbf{Z}$,
\item all the maps
$\mathcal{O}_Y(nd) \otimes_{\mathcal{O}_Y} \mathcal{O}_Y(m)
\longrightarrow
\mathcal{O}_Y(nd + m)$
of Equation (\ref{equation-multiply}) are isomorphisms,
\item all the maps
$\widetilde M(nd)|_Y = \widetilde M|_Y \otimes_{\mathcal{O}_Y}
\mathcal{O}_X(nd)|_Y \to \widetilde{M(nd)}|_Y$
(see \ref{equation-multiply-more-generally}) are isomorphisms,
\item given $f \in S_{nd}$ denote $s \in \Gamma(Y, \mathcal{O}_Y(nd))$
the image of $f$ via (\ref{equation-global-sections})
restricted to $Y$, then $D_{+}(f) \cap Y = Y_s$,
\item a basis for the topology on $Y$ is given
by the collection of opens $Y_s$, where $s \in \Gamma(Y, \mathcal{O}_Y(nd))$,
$n \geq 1$, and
\item a basis for the topology of $Y$ is given
by those opens $Y_s \subset Y$, for
$s \in \Gamma(Y, \mathcal{O}_Y(nd))$, $n \geq 1$ which are affine.
\end{enumerate}
\end{lemma}

\begin{proof}
Since $Y$ is quasi-compact there exist finitely many homogeneous
$f_i \in S_{+}$, $i = 1, \ldots, n$ such that the standard opens
$D_{+}(f_i)$ give an open covering of $Y$. Let $d_i = \deg(f_i)$ and set
$d = d_1 \ldots d_n$. Note that $D_{+}(f_i) = D_{+}(f_i^{d/d_i})$
and hence we see immediately that $Y \subset W_d$, by characterization
(2) in Lemma \ref{lemma-where-invertible} or
by (1) using Lemma \ref{lemma-when-invertible}.
Note that (1) implies (2), (3) and (4) by Lemma \ref{lemma-where-invertible}.
(Note that (3) is a special case of (4).)
Assertion (5) follows from Lemma \ref{lemma-principal-open}.
Assertions (6) and (7) follow because the open subsets $D_{+}(f)$
form a basis for the topology of $X$ and are affine.
\end{proof}

\begin{lemma}
\label{lemma-comparison-proj-quasi-coherent}
Let $S$ be a graded ring. Set $X = \text{Proj}(S)$.
Let $\mathcal{F}$ be a quasi-coherent $\mathcal{O}_X$-module.
Set $M = \bigoplus_{n \in \mathbf{Z}} \Gamma(X, \mathcal{F}(n))$ as
a graded $S$-module, using
(\ref{equation-global-sections-module}) and (\ref{equation-global-sections}).
Then there is a canonical $\mathcal{O}_X$-module map
$$
\widetilde{M} \longrightarrow \mathcal{F}
$$
functorial in $\mathcal{F}$ such that the induced map
$M_0 \to \Gamma(X, \mathcal{F})$ is the identity.
\end{lemma}

\begin{proof}
Let $f \in S$ be homogeneous of degree $d > 0$. Recall that
$\widetilde{M}|_{D_{+}(f)}$ corresponds to the
$S_{(f)}$-module $M_{(f)}$ by Lemma \ref{lemma-proj-sheaves}.
Thus we can define a canonical map
$$
M_{(f)} \longrightarrow \Gamma(D_+(f), \mathcal{F}),\quad
m/f^n \longmapsto m|_{D_+(f)} \otimes f|_{D_+(f)}^{-n}
$$
which makes sense because $f|_{D_+(f)}$ is a trivializing
section of the invertible sheaf $\mathcal{O}_X(d)|_{D_+(f)}$, see
Lemma \ref{lemma-when-invertible} and its proof.
Since $\widetilde{M}$ is quasi-coherent, this leads to a canonical
map
$$
\widetilde{M}|_{D_+(f)} \longrightarrow \mathcal{F}|_{D_+(f)}
$$
via Schemes, Lemma \ref{schemes-lemma-compare-constructions}.
We obtain a global map if we prove that the displayed maps glue on overlaps.
Proof of this is omitted. We also omit the proof of the final statement.
\end{proof}










